\documentclass[11pt]{article}

\usepackage[T1]{fontenc}
\usepackage[utf8]{inputenc}
\usepackage{lmodern}
\usepackage{geometry}
\usepackage{amsmath, amssymb}
\usepackage{graphicx}
\usepackage{booktabs}
\usepackage{hyperref}

\geometry{margin=1in}

\title{Project Title}
\author{
  Group number: XXX \\
  Student 1 \\
  Student 2 \\
  Student 3
}
\date{Modelling Complex Systems, Spring 2026}

\begin{document}

\maketitle

\section{Problem statement}

State clearly what model or problem you study and what questions you want to answer.

\section{Model and parameters}

Present the model equations, define the variables and parameters, and explain the modeling interpretation.

\section{Numerical method}

Describe the numerical procedure used in your code. Mention the main parameter values, the number of iterates, the discarded transient, and how the figures were produced.

\section{Results}

Present the main numerical results. Each figure should have a short caption and should be discussed in the text.

\begin{figure}[h!]
  \centering
  % \includegraphics[width=0.7\textwidth]{exampleFigure.pdf}
  \caption{Replace this caption with a description of your result.}
\end{figure}

\section{Interpretation}

Explain what the numerical results show. Connect the figures with the mathematical or modeling questions from the project.

\section{Conclusion}

Summarize the main conclusions in a few sentences.

\section{Use of LLM tools}

If you used an LLM in a meaningful way, briefly state:
\begin{itemize}
  \item which prompts were important,
  \item what output was useful,
  \item what had to be corrected or verified.
\end{itemize}

If you did not use an LLM, state that explicitly.

\section*{Files submitted}

List the files included in the submission.

\end{document}
